\section{Results, meaning, and comparison}
\label{sec:overview}

The central unrestricted-integer deliverable is a fully specified polynomial
\begin{equation}\label{eq:Q209}
 Q_{209}(\xi_0,\ldots,\xi_{208})
   =\sum_{j=0}^{148}R_j(\boldsymbol\xi)^2.
\end{equation}
Every \(R_j\) is an explicitly listed polynomial in \(\Z[\boldsymbol\xi]\)
of degree at most two. Appendix~\ref{app:residuals} prints every residual;
\artifact{quartic/equations/affine209_F64/equation.json} and its collected
expansion are the corresponding exact data. No parameter remains to be
chosen by an undecidable procedure.

\begin{theorem}[Consistency equivalence]\label{thm:main}
For the fixed polynomial \(Q_{209}\) in \eqref{eq:Q209},
\[
 \ZFC\vdash
 \left(\neg\exists\boldsymbol\xi\in\Z^{209}\;
                    Q_{209}(\boldsymbol\xi)=0\right)
       \longleftrightarrow\Con(\ZFC).
\]
The same statement holds for the delivered second quartic and the three
pooled polynomials, in their respective unrestricted-integer dimensions.
The separately specified polynomial \(P_{10}\) of
Section~\ref{sec:compact} satisfies
\[
 \ZFC\vdash
 \left(\neg\exists\boldsymbol x\in\Zp^{10}\;
                    P_{10}(\boldsymbol x)=0\right)
       \longleftrightarrow\Con(\ZFC).
\]
\end{theorem}

The proof composes the new arithmetic reductions below with the precise
inherited calibrations in Section~\ref{sec:boundary}. The consistency
predicate is the absence of a finite syntactic refutation in the
repository's fixed effective presentation \code{ZFCax_s}, with proof
relation \code{BProv} and contradiction \code{Form.fBot}. The quantifiers
are internal to ZFC; the assertion does not restrict proof codes to
externally standard integers. It assumes neither consistency nor
inconsistency and proves neither one.

\subsection{The delivered profiles}

\begin{center}
\small
\begin{tabular}{@{}lrrrrl@{}}
\toprule
Profile & Variables & Degree & Squares & Monomials & Height / cost\\
\midrule
Sparse quartic, \(\Z\) & 209 & 4 & 149 & 5,486 & height \(2^{32}\)\\
Lower-height quartic, \(\Z\) & 209 & 4 & 149 & 6,654 & height \(2^{30}\)\\
Pooled, \(\Z\) & 130 & 12 & 70 & circuit & 2,475 operations\\
Pooled, \(\Z\) & 116 & 16 & 56 & circuit & 2,433 operations\\
Pooled, \(\Z\) & 107 & 32 & 47 & circuit & 2,419 operations\\
Centered transport, \(\Zp\) & 10 & 190 & -- & circuit & 119 operations\\
\bottomrule
\end{tabular}
\end{center}

For the quartics, height is the largest absolute coefficient of the
\emph{collected final polynomial}, not a literal in its evaluator.
All monomial counts also refer to collected final polynomials. For the
circuit profiles, the finite straight-line program is the complete
polynomial specification; a full expanded monomial count is not claimed.
An operation is one live binary addition, subtraction, or multiplication;
squaring and multiplication by a fixed integer are charged, while reuse is
free. This is an arithmetic count, not a bit-complexity estimate.

The inspected repository snapshot is
\begin{center}\small
\code{d123610bc98fef2006bff1c9466bc3f1221003e7}.
\end{center}
It contains placed packages through \code{178-polynomial-centers}, together
with seven incoming ZIP archives. Comparisons below include those incoming
archives rather than treating the last placed package as the only baseline.

The immediate unrestricted quartic parent has 210 variables, 150 squared
quadratics, 5,282 monomials, and height \(2^{40}\). The sparse new quartic
saves one variable and one square and divides the height by 256, at the
cost of 204 additional monomials. The second new quartic divides the
parent's height by 1,024, at a larger support cost. These are explicit
tradeoffs, not simultaneous improvements in every measure.

The corresponding older pooled profiles have 131 variables at degree 12,
117 variables at degree 16, and 108 variables at degree 32. The new composition saves one variable
in each case, with unchanged degrees. For ten positive variables the
relevant points are source 172's degree 214 at 87 operations, source 178's
degree 208 at 116 operations, and the incoming affine-centering result's
degree 206 at 90 operations. Our degree 190 costs 119 operations. It
reduces the lowest degree among these inspected ten-positive-variable
constructions by 16, approximately 7.77\%, while using more operations.

There is no claim that 209 is the least possible integer dimension, that
190 is the least possible degree with ten positive variables, or that
these constructions are globally simplest. The repository already contains
much smaller-variable representations in other degree and coefficient
regimes. The word ``simplest'' requires a choice among dimension, degree,
coefficient height, support, evaluation length, and proof complexity.

\subsection{What is new in this iteration}

The first reduction replaces a separate positive square horizon by
\(t=2+4z\), and recovers its sign from the remaining exponent and endpoint
certificates. A large \emph{absolute} lower bound suffices before the sign
is known. This allows the sharp positive-center Pell bound to be applied
only after \(t>0\), avoiding an incorrect estimate on negative centers.
The retained divisibility then permits an integral rescaling of the growth
coordinate. A change of residual basis yields a second height tradeoff.

The compact reduction uses an integer polynomial approximation to an
entire Pell expression. The free quotient in its source admits a small
auxiliary index that invalidates a naive transport estimate. The new
transport \(L=UEb^2\) gives a uniform bound including that index, while
keeping the cross term below the degree of the centered norm defect.

The pooled reductions compose the corrected signed growth interface with
four-base reciprocal pooling and explicit acyclic eliminations. The pool
does not need pairwise coprime moduli. Its soundness comes from a strict
bound on a signed rational sum of squared errors.

\subsection{Why a different formal calculus alone is insufficient}

The direct branch uses the numerical calculus \(\Ccal\) through its
independent refutation calibration; it does not simulate the K215 machine
or instantiate Jones's universal polynomial. The compact branch retains
its native/compiler/K215 ancestry. These are different constructions.

Alternative logical encodings remain relevant, but they need a complete
unbounded compiler and explicit coefficients before yielding a competitive
equation. Formisano, Gavazzi, and Omodeo explore a relation-algebra
presentation of NBG in this direction~\cite{fgo}. General variable-count
theorems, including Sun's recent ten-integer-variable representation,
likewise do not by themselves specify a particularly small fixed
consistency instance~\cite{sun10}. The present report measures actual
emitted polynomials.
